\noindent\textit{Enunciado adaptado del problema 8 del Banco de Problemas de Aprendizaje por Refuerzo (Modelos de Decisión en el Tiempo, Módulo 3, IIND-3518, Universidad de los Andes).}

\bigskip

\noindent Un hotel escoge cada día una tarifa baja ($b$) o alta ($a$) según su ocupación: baja ($B$), media ($M$) o alta ($A$). El ingreso del día depende de la ocupación y de la tarifa:

\begin{table}[H]
\centering
\begin{tabular}{ccc}
\toprule
Ocupación & Tarifa baja ($b$) & Tarifa alta ($a$)\\
\midrule
B & 6,5 & 3\\
M & 6 & 4\\
A & 7 & 10\\
\bottomrule
\end{tabular}
\end{table}

\noindent La ocupación del día siguiente es \textbf{aleatoria}: desde una misma ocupación y con la misma tarifa, el hotel puede amanecer al día siguiente con ocupación baja, media o alta. Las probabilidades dependen solo de la ocupación de hoy y de la tarifa elegida, y están dadas por las siguientes tablas. Cada fila corresponde a la ocupación de hoy, cada columna a la ocupación de mañana, y cada fila suma 1.

\begin{table}[H]
\centering
\begin{tabular}{cccc}
\multicolumn{4}{c}{\textit{Si se cobra la tarifa baja ($b$)}}\\
\toprule
Ocupación hoy & Mañana B & Mañana M & Mañana A\\
\midrule
B & 0,3 & 0,5 & 0,2\\
M & 0,2 & 0,4 & 0,4\\
A & 0,1 & 0,3 & 0,6\\
\bottomrule
\end{tabular}
\qquad
\begin{tabular}{cccc}
\multicolumn{4}{c}{\textit{Si se cobra la tarifa alta ($a$)}}\\
\toprule
Ocupación hoy & Mañana B & Mañana M & Mañana A\\
\midrule
B & 0,7 & 0,3 & 0\\
M & 0,4 & 0,4 & 0,2\\
A & 0,2 & 0,5 & 0,3\\
\bottomrule
\end{tabular}
\end{table}

\noindent Por ejemplo, si hoy la ocupación es baja y se cobra la tarifa baja, mañana la ocupación será baja con probabilidad 0,3, media con probabilidad 0,5 y alta con probabilidad 0,2. Con ocupación baja y tarifa alta, en cambio, la ocupación no puede subir a alta al día siguiente. El algoritmo de Q-learning \textbf{no conoce estas probabilidades}: solo observa, en cada día simulado, la ocupación a la que efectivamente se pasa y el ingreso recibido.

\noindent El hotel opera indefinidamente, empieza con ocupación baja y usa $\gamma = 0{,}9$.

\begin{enumerate}
\item Formule el problema para Q-learning: estados, acciones factibles y qué observa el algoritmo después de cada decisión.
\item Entrene Q-learning con $\alpha = 0{,}01$, $\varepsilon = 0{,}2$, $2 \times 10^5$ pasos y semilla 2026. Presente la tabla $Q$ y la política greedy aprendida.
\item Calcule la política óptima de forma exacta con \emph{value iteration} y compárela con la aprendida en el inciso anterior. Compare también la ganancia máxima de cada estado con la ganancia de seguir la política aprendida.
\item Interprete la política óptima y compárela con la política greedy $(b, b, a)$ que se obtiene en el problema original después de procesar las seis experiencias.
\end{enumerate}

\noindent\textbf{Solución:} \texttt{qlearning\_problem\_solutions/qlearning\_hotel.ipynb}.
